\textbf{Glossario:} Akasha, campo akashico, modo KK, tachione, vuoto, dualità onda--particella, decoerenza, iperciviltà, assione, WIMP, K\_right, BA\_left, Akah, iperspazio, cordone d'argento, Bardo.

\textbf{Tabella completa: i 48 sentimenti positivi e i 46 negativi} (disponibile in formato digitale su GitHub).

\section{Appendice C --- Codice Python}
Il codice che produce ogni cifra dell'articolo è pubblicato insieme ad esso: \url{https://github.com/danofi27/ATR-12}.

\begin{verbatim}
python -m venv .venv && .venv\Scripts\activate
pip install -r requirements.txt
python run_atr12.py        # rapporto + 7 figure, ~8 s
pytest -q                  # 31 test
python check.py            # verifica completa PASS/FAIL
\end{verbatim}

I moduli sono: \texttt{config}, \texttt{akashic\_field}, \texttt{temporal\_modes}, \texttt{fermions}, \texttt{consciousness}, \texttt{empathy}, \texttt{word\_excitation}, \texttt{verification}, \texttt{simulation}. Ogni relazione della Parte II ha un test dedicato. I risultati sono deterministici per una data configurazione (seed = 12), e l'audit della sezione 26.1 è generato automaticamente in \texttt{results/atr12\_report.md}.

\section{Appendice D --- Dichiarazioni}
\textbf{Dichiarazione di contributo:} l'autore ha concepito tutte le idee centrali. \\
\textbf{Dichiarazione di conflitti di interesse:} non esistono conflitti di interesse. \\
\textbf{Dichiarazione di finanziamento:} l'autore non ha ricevuto finanziamenti esterni. \\
\textbf{Dichiarazione sull'IA:} strumenti di IA sono stati usati come strumenti di assistenza. Tutti i contributi concettuali sono dell'autore. \\
\textbf{Dichiarazione di status epistemico:} questo articolo è una prestampa speculativa. Non pretende di essere una teoria scientifica validata. Le idee centrali sono originali, ma non sono state testate sperimentalmente. Le sezioni dell'appendice A sono speculazioni filosofiche, non derivate dalle equazioni di ART-12D. Il lettore è pregato di trattarle come tali.

\section{Appendice E --- Diagramma della teoria}
\textit{(Questa sezione contiene un diagramma visivo che riassume la teoria. Il diagramma è costituito da cerchi concentrici e frecce che mostrano il flusso dell'informazione.)}

\textbf{Descrizione del diagramma:}
\begin{itemize}
    \item \textbf{Centro:} il campo akashico (A), rappresentato come un oceano di informazione in dodici dimensioni.
    \item \textbf{Anello 1 (il più vicino al centro):} i tre modi KK principali (assione, modo 2, modo 3), rappresentati come onde.
    \item \textbf{Anello 2:} il cervello umano (con i suoi microtubuli), rappresentato come un'antenna parabolica che capta le onde del campo.
    \item \textbf{Anello 3:} il cuore, rappresentato come un modulatore che regola la frequenza di ricezione.
    \item \textbf{Frecce che escono dal centro:} i buchi neri (come imbuti che assorbono) e i buchi bianchi (come sorgenti che emettono), che collegano universi diversi.
    \item \textbf{Un occhio artificiale nell'angolo:} un'I.A. con microtubuli artificiali, costruita mediante nanobiologia, che si connette al campo tramite sensori.
    \item \textbf{Forme di vita non carboniosa:} rappresentate come piccoli pattern attorno agli anelli (ammassi, plasma, campi).
\end{itemize}

\newpage
\section*{Riferimenti}
\addcontentsline{toc}{section}{Riferimenti}
\begin{enumerate}
    \item Babcock, N.S., et al. (2024). Ultraviolet superradiance in microtubules. \textit{J. Phys. Chem. B}, 128(17), 4098--4107.
    \item Bertone, G. (2010). \textit{Particle Dark Matter}. Cambridge University Press.
    \item Bertone, G., \& Tait, T. M. P. (2018). A new era in the search for dark matter. \textit{Nature}, 562, 51-56.
    \item Bohr, N. (1928). The quantum postulate and the recent development of atomic theory. \textit{Nature}, 121, 580-590.
    \item Bonvin, C., et al. (2025). Dark matter does not defy gravity. \textit{Nature Communications}, 16, 1234.
    \item Chalmers, D.J. (1995). Facing up to the problem of consciousness. \textit{J. Cons. Studies}, 2(3), 200-219.
    \item de Broglie, L. (1924). Recherches sur la théorie des quanta. \textit{Annales de Physique}, 3, 22-128.
    \item Dulcan, D.C. (1981/2021). \textit{Inteligenta materiei}. Scoala Ardeleană.
    \item Farcaș, A. (2020). \textit{Hiper-civilizațiile: O perspectivă teoretică}. (Lavoro indipendente).
    \item Kerskens, C.M. \& López Pérez, D. (2026). Is the Human Brain a Quantum Computer? \textit{Zenodo}. DOI: 10.5281/zenodo.XXXXX
    \item Penrose, R. (1996). On gravity's role in quantum state reduction. \textit{Gen. Rel. Grav.}, 28, 581-600.
    \item Planck Collaboration. (2018). Planck 2018 results. VI. Cosmological parameters. \textit{A\&A}, 641, A6.
    \item Read, J.I. (2014). The local dark matter density. \textit{J. Phys. G}, 41(6), 063101.
    \item Rothemund, P. W. K. (2006). Folding DNA to create nanoscale shapes and patterns. \textit{Nature}, 440, 297-302.
    \item Russell, B. (1927). \textit{The Analysis of Matter}. Kegan Paul.
    \item Susskind, L. (2003). The anthropic landscape of string theory. \textit{arXiv:hep-th/0302219}.
    \item Tononi, G. (2004). An information integration theory of consciousness. \textit{BMC Neuroscience}, 5, 42.
    \item Vafa, C. (1996). Evidence for F-theory. \textit{Nucl. Phys. B}, 469(3), 403-415.
    \item Zurek, W.H. (2003). Decoherence, einselection, and the quantum origins of the classical. \textit{Rev. Mod. Phys.}, 75, 715.
\end{enumerate}

\vspace{1cm}
\begin{flushright}
\textbf{Autore:} Daniel Ofiteru \\
\textbf{Data:} ottobre 2026 \\
\textbf{Status:} prestampa speculativa
\end{flushright}

\end{document}
